\documentclass[12pt,a4paper]{article}

% === ENCODAGE ET LANGUE ===
\usepackage[utf8]{inputenc}
\usepackage[T1]{fontenc}
\usepackage[french]{babel}

% === MISE EN PAGE ===
\usepackage{geometry}
\geometry{margin=2.5cm}
\usepackage{setspace}
\onehalfspacing

% === POLICES ET TYPOGRAPHIE ===
\usepackage{lmodern}
\usepackage{microtype}

% === GRAPHIQUES ET FIGURES ===
\usepackage{graphicx}
\graphicspath{{results/}}
\usepackage{tikz}
\usetikzlibrary{arrows.meta, shapes.geometric, positioning, calc, fit, backgrounds}
\usepackage{pgfplots}
\pgfplotsset{compat=1.18}
\usepackage{float}


% === TABLEAUX ===
\usepackage{booktabs}
\usepackage{multirow}
\usepackage{tabularx}
\usepackage{array}
\usepackage{colortbl}

% === COULEURS ===
\usepackage[table]{xcolor}
\definecolor{upsred}{RGB}{180, 30, 40}
\definecolor{criticalred}{RGB}{220, 53, 69}
\definecolor{importantyellow}{RGB}{255, 193, 7}
\definecolor{moderategray}{RGB}{108, 117, 125}
\definecolor{lightgray}{RGB}{245, 245, 245}
\definecolor{codebg}{RGB}{248, 249, 250}

% === MATHS ===
\usepackage{amsmath}
\usepackage{amssymb}

% === LISTINGS (CODE) ===
\usepackage{listings}
\lstset{
  basicstyle=\ttfamily\small,
  backgroundcolor=\color{codebg},
  frame=single,
  rulecolor=\color{gray!30},
  breaklines=true,
  columns=fullflexible,
  literate={é}{{\'e}}1 {è}{{\`e}}1 {ê}{{\^e}}1 {à}{{\`a}}1 {ù}{{\`u}}1 {û}{{\^u}}1 {ç}{{\c{c}}}1 {ô}{{\^o}}1 {î}{{\^i}}1
}

% === LIENS ===
\usepackage{hyperref}
\hypersetup{
  colorlinks=true,
  linkcolor=upsred,
  citecolor=upsred,
  urlcolor=blue!70!black,
  pdfauthor={Mohamed DIALLO et al.},
  pdftitle={GeoR2LLMs -- Rapport de travail Phase 1},
  pdfsubject={Système Question-Réponse Géographique avec RAG}
}

% === BIBLIOGRAPHIE ===
\usepackage{cite}

% === EN-TÊTES ET PIEDS DE PAGE ===
\usepackage{fancyhdr}
\pagestyle{fancy}
\fancyhf{}
\fancyhead[L]{\small\textit{GeoR2LLMs -- Rapport de travail Phase 1}}
\fancyhead[R]{\small\textit{M2 IA -- UPS}}
\fancyfoot[C]{\thepage}
\renewcommand{\headrulewidth}{0.4pt}
\renewcommand{\footrulewidth}{0.4pt}

\fancypagestyle{plain}{
  \fancyhf{}
  \fancyfoot[C]{\thepage}
  \renewcommand{\headrulewidth}{0pt}
  \renewcommand{\footrulewidth}{0pt}
}

% === TITRES ===
\usepackage{titlesec}
\titleformat{\section}{\Large\bfseries\color{upsred}}{\thesection}{1em}{}
\titleformat{\subsection}{\large\bfseries}{\thesubsection}{1em}{}
\titleformat{\subsubsection}{\normalsize\bfseries\itshape}{\thesubsubsection}{1em}{}

% === ENUMÉRATION ===
\usepackage{enumitem}

% === CAPTION ===
\usepackage{caption}
\captionsetup{font=small, labelfont=bf, format=hang}

% ====================================================================
\begin{document}

% ====================================================================
% PAGE DE GARDE
% ====================================================================
\begin{titlepage}
  \centering
  \vspace*{1cm}

  {\Large\textsc{Université Toulouse III -- Paul Sabatier}}\\[0.3cm]
  {\large\textsc{Master 2 Intelligence Artificielle}}\\[0.2cm]
  {\normalsize IRIT -- Institut de Recherche en Informatique de Toulouse}\\[2cm]

  \rule{\textwidth}{1.5pt}\\[0.6cm]
  {\huge\bfseries Système Question-Réponse\\[0.3cm] en Géographie avec RAG}\\[0.4cm]
  {\Large\textit{Rapport de travail -- Phase 1 : Plateforme de base}}\\[0.3cm]
  {\large Projet GeoR2LLMs}\\[0.4cm]
  \rule{\textwidth}{1.5pt}\\[2cm]

  \begin{minipage}{0.45\textwidth}
    \begin{flushleft}
      \textbf{Réalisé par :}\\[0.3cm]
      Mohamed DIALLO\\
      Hiba L'MOUDDEN\\
      Feras MAKHLOUF\\
      Osama BAAZZI\\
      Ousama BOUTEZROUT
    \end{flushleft}
  \end{minipage}
  \hfill
  \begin{minipage}{0.45\textwidth}
    \begin{flushright}
      \textbf{Sous la direction de :}\\[0.3cm]
      Madame Lynda TAMINE-LECHANI\\
      Professeure -- IRIT\\[0.5cm]
      \textbf{UE :} Chef d'\oe uvre\\
      \textbf{Année :} 2025--2026
    \end{flushright}
  \end{minipage}

  \vfill

  {\large Février 2026}
\end{titlepage}

% ====================================================================
% TABLE DES MATIÈRES
% ====================================================================
\newpage
\tableofcontents
\newpage

% ====================================================================
% 1. INTRODUCTION
% ====================================================================
\section{Introduction}
\label{sec:introduction}

Le projet GeoR2LLMs (\textit{Geographic Retrieval-augmented generation with Large Language Models}) s'inscrit dans le cadre du chef-d'\oe uvre du Master 2 Intelligence Artificielle de l'Université Toulouse III -- Paul Sabatier, sous la direction de Madame Lynda Tamine-Lechani à l'IRIT. Ce rapport de travail documente l'implémentation et les résultats de la Phase 1 du projet.

\subsection{Contexte et motivations}

Le Question-Réponse géographique (GeoQA) constitue un défi particulier pour les systèmes de traitement automatique des langues. Mai et al.~\cite{mai_geographic_2021} identifient les questions géographiques comme nécessitant un raisonnement spatial, la compréhension de relations topologiques et l'intégration de connaissances spécifiques au domaine. Cohn et Blackwell~\cite{cohn_evaluating_2024} montrent qu'aucun LLM actuel ne parvient à raisonner de manière fiable sur les directions cardinales, même avec une température de zéro.

L'architecture RAG (\textit{Retrieval-Augmented Generation}), introduite par Lewis et al.~\cite{lewis2020retrieval} dans le cadre de NeurIPS 2020, propose d'augmenter les LLMs avec des connaissances externes récupérées dynamiquement depuis un corpus. Gao et al.~\cite{gao2023ragsurvey} proposent une taxonomie des systèmes RAG distinguant le RAG na\"if (retrieve-then-read) du RAG avancé (avec réécriture de requêtes ou re-ranking). Notre implémentation suit l'approche na\"ive.

\subsection{Objectif de la Phase 1}

La Phase 1 du projet vise à mettre en place une \textbf{plateforme de base} pour le système GeoQA-RAG, conformément au cahier des charges. Les objectifs sont :
\begin{enumerate}[leftmargin=*]
  \item Implémenter un pipeline RAG complet avec les composants spécifiés (Qwen3-0.6B, BGE-m3, FAISS)
  \item Établir une baseline (LLM seul, sans RAG) comme point de comparaison
  \item Évaluer l'apport du RAG sur le dataset GeoSQA~\cite{huang_geosqa_2019}
  \item Mesurer la qualité du retrieval indépendamment de la génération
\end{enumerate}

\subsection{Organisation du rapport}

La section~\ref{sec:architecture} présente l'architecture du système. La section~\ref{sec:methodologie} détaille les choix méthodologiques. Les résultats expérimentaux sont rapportés en section~\ref{sec:resultats} et analysés en section~\ref{sec:analyse}. La section~\ref{sec:conclusion} conclut et ouvre sur la Phase 2.

% ====================================================================
% 2. ARCHITECTURE DU SYSTÈME
% ====================================================================
\section{Architecture du système}
\label{sec:architecture}

Le système repose sur une architecture RAG composée de trois modules : un \textbf{Retriever}, un \textbf{Generator} et un \textbf{Evaluator}. La figure~\ref{fig:architecture} illustre le pipeline.

\begin{figure}[H]
\centering
\resizebox{\textwidth}{!}{%
\begin{tikzpicture}[
  node distance=1.2cm and 2cm,
  box/.style={rectangle, draw, rounded corners=4pt, minimum height=1.2cm, minimum width=3cm, align=center, font=\small},
  data/.style={box, fill=blue!8, draw=blue!40},
  process/.style={box, fill=green!8, draw=green!40},
  model/.style={box, fill=orange!10, draw=orange!40},
  result/.style={box, fill=red!8, draw=red!40},
  arrow/.style={-{Stealth[length=3mm]}, thick, draw=gray!70},
  lbl/.style={font=\scriptsize\itshape, text=gray}
]

% Données d'entrée
\node[data] (wiki) {Wikipedia FR\\{\scriptsize 300 articles}};
\node[data, below=0.8cm of wiki] (question) {Question\\{\scriptsize GeoSQA}};

% Retriever
\node[process, right=2cm of wiki] (chunking) {Chunking\\{\scriptsize 512 car., overlap 50}};
\node[model, right=1.5cm of chunking] (bge) {BGE-m3\\{\scriptsize 1024 dim.}};
\node[process, right=1.5cm of bge] (faiss) {FAISS\\{\scriptsize IndexFlatIP}};

% Generator
\node[model, right=2cm of question] (embed_q) {Encodage\\{\scriptsize question}};
\node[process, right=3.7cm of embed_q] (topk) {Top-K\\{\scriptsize K=5}};
\node[model, below=2cm of faiss] (qwen) {Qwen3-0.6B\\{\scriptsize Generator}};
\node[result, right=1.5cm of qwen] (reponse) {Réponse\\{\scriptsize générée}};

% Flèches corpus
\draw[arrow] (wiki) -- (chunking);
\draw[arrow] (chunking) -- (bge);
\draw[arrow] (bge) -- node[above, lbl] {embeddings} (faiss);

% Flèches question
\draw[arrow] (question) -- (embed_q);
\draw[arrow] (embed_q) -- node[above, lbl] {vecteur requête} (topk);
\draw[arrow] (faiss) -- node[right, lbl] {index} (topk);

% Flèches génération
\draw[arrow] (topk) -- node[right, lbl] {contexte} (qwen);
\draw[arrow] (question.south) |- (qwen.west);
\draw[arrow] (qwen) -- (reponse);

% Cadres
\begin{scope}[on background layer]
  \node[fit=(chunking)(bge)(faiss), draw=green!60, dashed, rounded corners=8pt, inner sep=10pt, label={[font=\small\bfseries, text=green!60!black]above:Module Retriever}] {};
  \node[fit=(qwen)(reponse), draw=orange!60, dashed, rounded corners=8pt, inner sep=10pt, label={[font=\small\bfseries, text=orange!60!black]below:Module Generator}] {};
\end{scope}

\end{tikzpicture}
}
\caption{Architecture du pipeline RAG. Le corpus Wikipedia est découpé en chunks, encodé par BGE-m3 et indexé dans FAISS. Les questions sont encodées puis les top-K passages récupérés sont fournis comme contexte au générateur Qwen3-0.6B.}
\label{fig:architecture}
\end{figure}

\subsection{Module Retriever}

Le retriever utilise le modèle d'embeddings multilingue \textbf{BGE-m3} (BAAI), présenté par Chen et al.~\cite{chen2024bgem3} à ACL 2024 Findings. BGE-m3 produit des vecteurs denses de 1024 dimensions et supporte plus de 100 langues. Les documents sont encodés et stockés dans un index \textbf{FAISS}~\cite{johnson2019faiss} (\texttt{IndexFlatIP}) utilisant le produit scalaire après normalisation L2, équivalent à la similarité cosinus. Pour une requête $q$, le retriever retourne les $K$ passages les plus similaires :
\begin{equation}
  \text{Top-K}(q) = \underset{d_i \in \mathcal{C}}{\text{argmax-K}} \; \cos(\mathbf{e}_q, \mathbf{e}_{d_i})
\end{equation}
où $\mathbf{e}_q$ et $\mathbf{e}_{d_i}$ sont les embeddings de la requête et du document, et $\mathcal{C}$ est le corpus de chunks.

\subsection{Module Generator}

Le générateur est basé sur \textbf{Qwen3-0.6B}~\cite{yang2025qwen3}, un modèle de langage de 0.6 milliard de paramètres de la famille Qwen3 (Alibaba). Selon le rapport technique Qwen3, ce modèle supporte un mode \textit{thinking} activable. Nous l'utilisons avec \texttt{enable\_thinking=False} pour obtenir des réponses directes. Le prompt est structuré via le template de chat Qwen3 :

\begin{lstlisting}[language=Python, caption={Prompt template pour le mode RAG}]
system: "Tu es un expert en geographie. Reponds
         uniquement avec la lettre correcte."
user:   "Contexte: {documents recuperes}
         Question: {question}
         Options: A) ... B) ... C) ... D) ..."
\end{lstlisting}

En mode \textbf{baseline}, aucun contexte n'est fourni : le modèle répond uniquement à partir de ses connaissances internes.

\subsection{Module Evaluator}

L'évaluateur mesure la performance selon deux axes, conformément aux métriques standard du QA~\cite{rajpurkar2016squademnlp} :
\begin{itemize}[leftmargin=*]
  \item \textbf{Qualité de génération} : Accuracy MCQ, Exact Match, F1 token-level, Contains
  \item \textbf{Qualité de retrieval} : Hit Rate@K, MRR@K (métriques classiques en recherche d'information~\cite{karpukhin2020dense})
\end{itemize}

% ====================================================================
% 3. MÉTHODOLOGIE
% ====================================================================
\section{Méthodologie}
\label{sec:methodologie}

\subsection{Configuration expérimentale}

Le tableau~\ref{tab:config} résume l'ensemble des hyperparamètres et composants du système, tels qu'enregistrés dans le fichier \texttt{phase1\_metrics.json} produit par le notebook.

\begin{table}[H]
\centering
\caption{Configuration du système GeoR2LLMs -- Phase 1}
\label{tab:config}
\begin{tabular}{@{} l l l @{}}
\toprule
\textbf{Catégorie} & \textbf{Paramètre} & \textbf{Valeur} \\
\midrule
\multirow{5}{*}{\textbf{Génération}}
  & Modèle LLM & Qwen/Qwen3-0.6B \\
  & Temperature & 0.7 \\
  & Top-p (nucleus) & 0.8 \\
  & Top-k (sampling) & 20 \\
  & Max tokens générés & 100 \\
\midrule
\multirow{4}{*}{\textbf{Retrieval}}
  & Modèle embeddings & BAAI/bge-m3 (1024 dim.) \\
  & Index vectoriel & FAISS IndexFlatIP \\
  & Taille des chunks & 512 caractères \\
  & Chevauchement (overlap) & 50 caractères \\
\midrule
\multirow{2}{*}{\textbf{Dataset}}
  & Benchmark & rfr2003/Geo\_Benchmark \\
  & Configuration & GeoSQA (test split) \\
\midrule
\multirow{3}{*}{\textbf{Corpus}}
  & Source & wikimedia/wikipedia \\
  & Version & 20231101.fr \\
  & Nombre d'articles & 300 \\
\midrule
\textbf{Environnement} & GPU & Google Colab -- NVIDIA T4 (16 Go) \\
\bottomrule
\end{tabular}
\end{table}

\subsection{Dataset : GeoSQA}

GeoSQA~\cite{huang_geosqa_2019} est un benchmark de question-réponse scénarisé en géographie, développé par Huang et al. à l'Université de Nanjing et présenté à EMNLP-IJCNLP 2019. Le dataset contient \textbf{1\,981 scénarios} et \textbf{4\,110 questions à choix multiples} au niveau lycée chinois, où les diagrammes ont été annotés en descriptions textuelles. Chaque question comprend :
\begin{itemize}[leftmargin=*]
  \item Un \textbf{scénario} contextuel (description de situation géographique)
  \item La \textbf{question} associée
  \item \textbf{4 options} (A, B, C, D)
  \item La \textbf{réponse correcte} (une lettre)
\end{itemize}

Le split de test contient \textbf{838 questions}. L'évaluation porte sur un sous-ensemble de \textbf{200 questions} échantillonnées aléatoirement, pour des raisons de temps d'exécution sur Google Colab.

\paragraph{Langue des questions.} Les questions GeoSQA sont originellement en chinois et ont été traduites en anglais dans la version HuggingFace (\texttt{rfr2003/Geo\_Benchmark}). Notre corpus est constitué d'articles Wikipédia en français, créant un décalage linguistique entre les questions et le corpus.

\subsection{Corpus : Wikipedia français}

Le corpus est constitué de \textbf{300 articles} chargés en streaming depuis le dump Wikipedia français (\texttt{wikimedia/wikipedia}, configuration \texttt{20231101.fr}). Chaque article est découpé en chunks de 512 caractères avec un chevauchement de 50 caractères, produisant un total de \textbf{2\,527 chunks} (valeur mesurée lors de l'exécution).

\subsection{Protocole d'évaluation}

\subsubsection{Métriques de génération}

Les métriques de génération suivent les conventions établies par Rajpurkar et al.~\cite{rajpurkar2016squademnlp} pour le benchmark SQuAD :

\begin{itemize}[leftmargin=*]
  \item \textbf{Accuracy (MCQ)} : Proportion de questions où la lettre prédite correspond à la lettre correcte.
  \begin{equation}
    \text{Accuracy} = \frac{1}{N} \sum_{i=1}^{N} \mathbb{1}[\hat{y}_i = y_i]
  \end{equation}

  \item \textbf{Exact Match (EM)} : Correspondance exacte entre le texte prédit (normalisé) et le texte de la réponse correcte.

  \item \textbf{F1 Score} : F1 au niveau des tokens entre la prédiction et la référence, après tokenisation et normalisation.
  \begin{equation}
    F_1 = \frac{2 \cdot P \cdot R}{P + R}, \quad P = \frac{|\hat{T} \cap T|}{|\hat{T}|}, \quad R = \frac{|\hat{T} \cap T|}{|T|}
  \end{equation}
  où $\hat{T}$ et $T$ sont les ensembles de tokens de la prédiction et de la référence.

  \item \textbf{Contains} : Proportion de cas où les mots-clés de la réponse correcte apparaissent dans la prédiction.
\end{itemize}

\subsubsection{Métriques de retrieval}

Les métriques de retrieval suivent les conventions du \textit{Dense Passage Retrieval} (DPR) de Karpukhin et al.~\cite{karpukhin2020dense} :

\begin{itemize}[leftmargin=*]
  \item \textbf{Hit Rate@K} : Proportion de questions pour lesquelles au moins un des $K$ documents récupérés contient le texte de la réponse correcte.
  \begin{equation}
    \text{HR@K} = \frac{1}{N} \sum_{i=1}^{N} \mathbb{1}\left[\exists j \leq K : y_i \in d_j^{(i)}\right]
  \end{equation}

  \item \textbf{MRR@K} (\textit{Mean Reciprocal Rank}) : Moyenne du rang réciproque du premier document pertinent parmi les top-K.
  \begin{equation}
    \text{MRR@K} = \frac{1}{N} \sum_{i=1}^{N} \frac{1}{\text{rank}_i}
  \end{equation}
\end{itemize}

% réponse à : Pourquoi exact match est à 0, comment calcluez vous accuracy et exact match.
\subsubsection{Distinction Accuracy MCQ et Exact Match : pourquoi EM = 0\,\%}
\label{sec:em_clarification}

Il est essentiel de distinguer deux métriques souvent confondues dans ce contexte :

\begin{itemize}[leftmargin=*]
  \item \textbf{Accuracy MCQ} : cette métrique extrait par expression régulière la
  première lettre valide (A, B, C ou D) présente dans la réponse générée, puis la
  compare à la lettre de référence. Elle est \emph{robuste à la verbosité} : même
  si le modèle génère une phrase complète, la lettre est extraite correctement.
  C'est la métrique \textbf{principale} pour évaluer la compréhension des QCM.

  \item \textbf{Exact Match (EM)} : cette métrique compare le \emph{texte complet
  normalisé} de la prédiction avec le texte complet de la réponse de référence.
  Dans GeoSQA, la référence est une lettre seule (ex. \texttt{"B"}), alors que le
  modèle génère systématiquement une phrase complète, parfois dans une langue
  différente. Cette divergence de format rend l'EM \textbf{structurellement nul},
  indépendamment de la qualité de la réponse.
\end{itemize}

\begin{figure}[H]
\centering
\begin{tikzpicture}[
  box/.style={rectangle, draw, rounded corners=4pt, minimum height=1.2cm,
              align=left, font=\small, text width=9cm, inner sep=8pt},
  ok/.style={draw=green!60!black, fill=green!5},
  bad/.style={draw=red!60, fill=red!5},
  neutral/.style={draw=gray!50, fill=gray!5},
  arrow/.style={->, thick, color=gray!70}
]
  \node[neutral, box] (ref) at (0,0)
    {\textbf{Référence GeoSQA :}\\[4pt]
     \texttt{"B"}};

  \node[neutral, box, below=0.6cm of ref] (pred)
    {\textbf{Prédiction Qwen3-0.6B :}\\[4pt]
     \texttt{"La réponse correcte est B) La production de légumes}\\
     \texttt{concentrée à Shouguang..."}};

  \node[ok, box, below=0.6cm of pred] (acc)
    {\textcolor{green!50!black}{\checkmark}~\textbf{Accuracy MCQ = 1}\\[4pt]
     (extraction regex : lettre \texttt{"B"} trouvée = lettre de référence)};

  \node[bad, box, below=0.6cm of acc] (em)
    {\textcolor{red}{\texttimes}~\textbf{Exact Match = 0}\\[4pt]
     (texte complet prédit $\neq$ \texttt{"B"} après normalisation)};

  \draw[arrow] (ref.south) -- ++(0,-0.1) -| ([xshift=-1.2cm]acc.west) -- (acc.west);
  \draw[arrow] (pred.south) -- ++(0,-0.1) -| ([xshift=1.2cm]em.east) -- (em.east);
\end{tikzpicture}
\caption{Illustration de la divergence entre Accuracy MCQ et Exact Match.
  L'EM est structurellement nul car le modèle génère des réponses verboses
  alors que la référence est une lettre seule.
  L'Accuracy MCQ reste la métrique fiable dans ce contexte MCQ.}
\label{fig:em_explanation}
\end{figure}

\noindent\textbf{Conséquence pratique.} Dans la suite de ce rapport, l'Exact Match
est conservé pour des raisons de comparabilité avec la
littérature~\cite{rajpurkar2016squademnlp}, mais \textbf{l'Accuracy MCQ est la
métrique de référence} pour juger de la performance du système sur GeoSQA.



% ====================================================================
% 4. RÉSULTATS EXPÉRIMENTAUX
% ====================================================================
\section{Résultats expérimentaux}
\label{sec:resultats}

Tous les résultats présentés dans cette section sont issus de l'exécution du notebook \texttt{Phase1\_GeoR2LLMs.ipynb} sur Google Colab (GPU T4). Les valeurs numériques sont extraites du fichier \texttt{phase1\_metrics.json} et les prédictions détaillées du fichier \texttt{phase1\_predictions.csv}, tous deux générés automatiquement par le notebook.

\subsection{Stabilité du modèle face aux variations de prompt}
\label{sec:prompt_stability}

Avant de comparer Baseline et RAG, nous avons testé \textbf{5 stratégies de
formulation} sur 50 questions GeoSQA afin de vérifier que les performances mesurées
ne dépendent pas du prompt choisi, conformément aux recommandations de
Brown et al.~\cite{brown2020language} sur l'impact du \textit{prompt engineering}.

\begin{table}[H]
\centering
\caption{Résultats du test de stabilité sur 5 stratégies de prompt (50 questions).
  Classement par accuracy décroissante. Le seuil aléatoire est de 25\,\%.}
\label{tab:prompt_stability}
\begin{tabular}{@{} c l c c @{}}
\toprule
\textbf{Rang} & \textbf{Stratégie} & \textbf{Accuracy MCQ} & \textbf{Écart vs P1} \\
\midrule
1 & P1 — Prompt basique (FR)        & \textbf{52.0\,\%} & --      \\
2 & P4 — Few-Shot 2 exemples (FR)   & 38.0\,\%          & --14 pts \\
2 & P5 — Prompt basique (EN)        & 38.0\,\%          & --14 pts \\
4 & P2 — Expert contraint (FR)      & 30.0\,\%          & --22 pts \\
5 & P3 — Chain-of-Thought (FR)      & 20.0\,\%          & --32 pts \\
\midrule
\multicolumn{2}{@{}l}{\textit{Seuil aléatoire (4 choix)}} & 25.0\,\% & \\
\multicolumn{2}{@{}l}{\textit{Écart max (P1 -- P3)}}      & \textbf{32 pts}  & \\
\bottomrule
\end{tabular}
\end{table}

\begin{figure}[H]
\centering
\begin{tikzpicture}
\begin{axis}[
  xbar,
  width=0.82\textwidth,
  height=6.2cm,
  bar width=16pt,
  % On définit les étiquettes sans espaces inutiles autour des accolades
  symbolic y coords={P3 Chain-of-Thought,P2 Expert contraint,P5 Basique EN,P4 Few-Shot,P1 Basique FR},
  ytick=data,
  yticklabel style={font=\small},
  xmin=0, xmax=62,
  xtick={0,10,20,25,30,40,50,60},
  xlabel={Accuracy MCQ (\%)},
  xlabel style={font=\small},
  nodes near coords,
  nodes near coords style={font=\scriptsize, /pgf/number format/.cd, fixed, precision=0},
  every node near coord/.append style={xshift=4pt},
  nodes near coords align={horizontal},
  title={Accuracy MCQ par stratégie de prompt (50 questions GeoSQA)},
  title style={font=\small\bfseries},
  ymajorgrids=true,
  grid style={dashed, gray!30},
  legend style={at={(0.98,0.05)}, anchor=south east, font=\small},
]

\addplot[fill=blue!30, draw=blue!60] coordinates {
  (52,P1 Basique FR)
  (38,P4 Few-Shot)
  (38,P5 Basique EN)
  (30,P2 Expert contraint)
  (20,P3 Chain-of-Thought)
};

\addplot[dashed, color=red!70, sharp plot, no markers, line width=1.3pt]
  coordinates { (25,P3 Chain-of-Thought) (25,P1 Basique FR) };

\addlegendentry{Accuracy MCQ}
\addlegendentry{Seuil aléatoire (25\%)}
\end{axis}
\end{tikzpicture}
\caption{Accuracy MCQ selon la stratégie de prompt (50 questions GeoSQA). La ligne rouge pointillée indique le seuil aléatoire (25\,\%). L'écart maximal de 32 points entre P1 et P3 montre que la formulation influence fortement un modèle de 0.6B, contrairement aux grands modèles. Le prompt P3 Chain-of-Thought dégrade même les performances sous le seuil aléatoire, suggérant que le raisonnement intermédiaire perturbe l'extraction de la lettre finale.}
\label{fig:prompt_stability}
\end{figure}

\paragraph{Choix du prompt de référence.}
Contrairement aux grands modèles où les stratégies avancées (CoT, Few-Shot)
apportent des gains significatifs~\cite{wei_chain--thought_2022}, Qwen3-0.6B
est très sensible à la complexité de la consigne. Le prompt P1 (basique FR)
obtient la meilleure accuracy (52\,\% sur 50 questions), mais cet écart avec
les évaluations sur 200 questions (24.5\,\%) indique une variance élevée sur
petit échantillon. Pour les évaluations principales sur 200 questions, nous
utilisons le \textbf{prompt P2 (Expert contraint)} car il impose le format de
sortie le plus strict (lettre seule), réduisant les erreurs d'extraction regex
et garantissant la comparabilité avec la Phase~2.


\subsection{Comparaison Baseline vs RAG}
\label{sec:baseline_vs_rag}

Le tableau~\ref{tab:results_main} présente les résultats complets de l'évaluation
sur 200 questions GeoSQA, avec les \textbf{taux d'accroissement relatifs}
$\Delta_r = (v_{\text{RAG}} - v_{\text{Base}}) / v_{\text{Base}} \times 100$.

\begin{table}[H]
\centering
\caption{Résultats Baseline vs RAG sur GeoSQA (200 questions).
  $\Delta_{abs}$ : différence absolue.
  $\Delta_{rel}$ : taux d'accroissement relatif
  $= (v_{\text{RAG}} - v_{\text{Base}}) / v_{\text{Base}} \times 100$.
  Valeurs issues de \texttt{phase1\_metrics.json}.}
\label{tab:results_main}
\begin{tabular}{@{} l c c c c @{}}
\toprule
\textbf{Métrique} & \textbf{Baseline} & \textbf{RAG}
  & \textbf{$\Delta_{abs}$} & \textbf{$\Delta_{rel}$ (\%)} \\
\midrule
Accuracy (MCQ) & 0.2450 & 0.2700
  & \textcolor{green!60!black}{+0.0250}
  & \textcolor{green!60!black}{+10.2\,\%} \\
Exact Match    & 0.0000 & 0.0000
  & 0.0000 & -- \\
F1 Score       & 0.0299 & 0.0114
  & \textcolor{red}{--0.0185}
  & \textcolor{red}{--61.9\,\%} \\
Contains       & 0.0300 & 0.0250
  & \textcolor{red}{--0.0050}
  & \textcolor{red}{--16.7\,\%} \\
\midrule
Temps (s)      & 348.5  & 486.0
  & +137.5 & \textcolor{orange}{+39.5\,\%} \\
\bottomrule
\end{tabular}
\end{table}

\begin{figure}[H]
\centering
\begin{tikzpicture}
\begin{axis}[
  ybar=4pt,
  width=0.92\textwidth,
  height=7cm,
  bar width=18pt,
  symbolic x coords={Accuracy (MCQ), Exact Match, F1 Score, Contains},
  xtick=data,
  xticklabel style={font=\small},
  ymin=0, ymax=0.35, 
  ytick={0, 0.05, 0.1, 0.15, 0.2, 0.25, 0.3},
  yticklabel style={font=\scriptsize, /pgf/number format/.cd, fixed, precision=2},
  ylabel={Score},
  ylabel style={font=\small},
  ymajorgrids=true,
  grid style={dashed, gray!30},
  nodes near coords,
  % --- CORRECTION ICI ---
  nodes near coords style={
    font=\tiny, 
    yshift=2pt, % Sorti du bloc /pgf/number format/
    /pgf/number format/.cd, fixed, precision=3
  },
  % -----------------------
  title={Comparaison Baseline vs RAG --- Métriques de Génération},
  title style={font=\small\bfseries},
  legend style={at={(0.02,0.95)}, anchor=north west, font=\tiny},
  legend columns=2,
  clip=true,
]

\addplot[fill=blue!35, draw=blue!70] coordinates {
  (Accuracy (MCQ), 0.2450)
  (Exact Match,    0.0000)
  (F1 Score,       0.0299)
  (Contains,       0.0300)
};

\addplot[fill=red!35, draw=red!70] coordinates {
  (Accuracy (MCQ), 0.2700)
  (Exact Match,    0.0000)
  (F1 Score,       0.0114)
  (Contains,       0.0250)
};

\addlegendentry{Baseline (LLM seul)}
\addlegendentry{RAG (LLM + Wikipedia)}
\end{axis}
\end{tikzpicture}
\caption{Comparaison des métriques de génération (200 questions GeoSQA).Le RAG améliore l'Accuracy MCQ de +10.2\,\% en relatif (+0.025 en absolu),mais dégrade le F1 de --61.9\,\% et le Contains de --16.7\,\%.L'Exact Match est nul dans les deux cas (cf. section~\ref{sec:em_clarification}).}
\label{fig:baseline_vs_rag}
\end{figure}

\paragraph{Distribution des prédictions Baseline.}
L'examen des 200 prédictions Baseline révèle un biais massif vers la lettre A :
123 réponses sur 200 (61.5\,\%) sont \texttt{A}, contre seulement 8 (4.0\,\%)
pour \texttt{B} et 10 (5.0\,\%) pour \texttt{C}. Ce biais confirme que Qwen3-0.6B,
faute de connaissances géographiques suffisantes, adopte une stratégie de repli
systématique vers le premier choix.

\begin{table}[H]
\centering
\caption{Distribution des lettres prédites en mode Baseline (200 questions).
  La distribution attendue sous H$_0$ (réponse aléatoire uniforme) serait
  50 occurrences par lettre (25\,\%).}
\label{tab:distribution_baseline}
\begin{tabular}{@{} l c c c @{}}
\toprule
\textbf{Lettre} & \textbf{Prédictions} & \textbf{\% observé}
  & \textbf{\% attendu (H$_0$)} \\
\midrule
A & 123 & \textbf{61.5\,\%} & 25.0\,\% \\
B & 8   & 4.0\,\%            & 25.0\,\% \\
C & 10  & 5.0\,\%            & 25.0\,\% \\
D & 59  & 29.5\,\%           & 25.0\,\% \\
\midrule
\textbf{Total} & \textbf{200} & \textbf{100\,\%} & 100\,\% \\
\bottomrule
\end{tabular}
\end{table}

\begin{figure}[H]
\centering
\begin{tikzpicture}
\begin{axis}[
  ybar=4pt,
  width=0.72\textwidth,
  height=5.5cm,
  bar width=22pt,
  symbolic x coords={A, B, C, D},
  xtick=data,
  xticklabel style={font=\normalsize\bfseries},
  ymin=0, ymax=75,
  ytick={0, 10, 20, 25, 30, 40, 50, 60, 70},
  ylabel={Nombre de prédictions},
  ylabel style={font=\small},
  ymajorgrids=true,
  grid style={dashed, gray!30},
  nodes near coords,
  nodes near coords style={font=\scriptsize},
  title={Distribution des prédictions Baseline par lettre},
  title style={font=\small\bfseries},
  legend style={at={(0.98,0.98)}, anchor=north east, font=\small},
]
% Prédictions réelles
\addplot[fill=blue!30, draw=blue!60] coordinates {
  (A, 123) (B, 8) (C, 10) (D, 59)
};
% Distribution uniforme attendue
\addplot[fill=gray!25, draw=gray!50, dashed] coordinates {
  (A, 50) (B, 50) (C, 50) (D, 50)
};
\addlegendentry{Prédictions Baseline}
\addlegendentry{Distribution uniforme (H$_0$)}
\end{axis}
\end{tikzpicture}
\caption{Distribution des lettres prédites par Qwen3-0.6B en mode Baseline
  (200 questions). Le modèle prédit \texttt{A} dans 61.5\,\% des cas, révélant
  un biais de position (position bias) caractéristique des modèles de petite
  taille manquant de connaissances factuelles. La distribution uniforme attendue
  sous l'hypothèse de réponse aléatoire (H$_0$) est indiquée en gris.}
\label{fig:distribution_baseline}
\end{figure}


\subsection{Métriques de retrieval}

Le tableau~\ref{tab:retrieval} et la figure~\ref{fig:retrieval} présentent les performances du module de retrieval.

\begin{table}[H]
\centering
\caption{Métriques de retrieval (BGE-m3 + FAISS). Valeurs issues de \texttt{phase1\_metrics.json}.}
\label{tab:retrieval}
\begin{tabular}{@{} l c @{}}
\toprule
\textbf{Métrique} & \textbf{Score} \\
\midrule
Hit Rate@1 & 0.040 \\
Hit Rate@3 & 0.060 \\
Hit Rate@5 & 0.080 \\
MRR@5 & 0.053 \\
\bottomrule
\end{tabular}
\end{table}

\begin{figure}[H]
\centering
\includegraphics[width=0.75\textwidth]{fig2_retrieval_metrics.png}
\caption{Métriques de retrieval (Hit Rate@K et MRR@5). Graphique généré par le notebook.}
\label{fig:retrieval}
\end{figure}

\paragraph{Scores de similarité cosinus.}
Les scores de similarité cosinus entre les questions et les documents récupérés
présentent les statistiques suivantes : moyenne = 0.477, médiane = 0.474,
minimum = 0.334, maximum = 0.652. Ces valeurs, supérieures à 0.5 pour les
meilleurs résultats, indiquent que BGE-m3 trouve des documents
\emph{superficiellement} similaires, mais cette similarité sémantique ne suffit
pas à garantir la pertinence factuelle. La similarité inter-lingue
(question EN, documents FR) plafonne la qualité du retrieval indépendamment
du modèle d'embedding utilisé.

\subsection{Impact de la valeur de $K$ sur les performances RAG}
\label{sec:impact_k}

Nous avons évalué $K \in \{1, 3, 5, 10\}$ sur 30 questions GeoSQA afin d'analyser
le compromis entre apport d'information et saturation du contexte.

\begin{table}[H]
\centering
\caption{Accuracy MCQ RAG en fonction de $K$ (30 questions GeoSQA).
  $\Delta_{abs}$ et $\Delta_{rel}$ calculés par rapport à $K=1$.}
\label{tab:impact_k}
\begin{tabular}{@{} c c c c @{}}
\toprule
\textbf{$K$} & \textbf{Accuracy MCQ}
  & \textbf{$\Delta_{abs}$ vs $K=1$}
  & \textbf{$\Delta_{rel}$ vs $K=1$ (\%)} \\
\midrule
1  & 23.3\,\% & --               & --         \\
3  & \textbf{30.0\,\%}
   & \textcolor{green!60!black}{+6.7 pts}
   & \textcolor{green!60!black}{+28.6\,\%}  \\
5  & 26.7\,\%
   & \textcolor{green!60!black}{+3.3 pts}
   & \textcolor{green!60!black}{+14.3\,\%}  \\
10 & \textbf{30.0\,\%}
   & \textcolor{green!60!black}{+6.7 pts}
   & \textcolor{green!60!black}{+28.6\,\%}  \\
\midrule
Baseline (sans RAG) & 24.5\,\% & \multicolumn{2}{l}{(référence générale)} \\
\bottomrule
\end{tabular}
\end{table}

\begin{figure}[H]
\centering
\begin{tikzpicture}
\begin{axis}[
  width=0.80\textwidth,
  height=6cm,
  xlabel={Valeur de $K$ (nombre de documents récupérés)},
  ylabel={Accuracy MCQ (\%)},
  xlabel style={font=\small},
  ylabel style={font=\small},
  xmin=0, xmax=11,
  ymin=15, ymax=38,
  xtick={1, 3, 5, 10},
  ytick={15, 20, 23.3, 24.5, 25, 26.7, 30, 35},
  yticklabel style={font=\scriptsize,
    /pgf/number format/.cd, fixed, precision=1},
  ymajorgrids=true,
  grid style={dashed, gray!30},
  mark size=3.5pt,
  title={Accuracy RAG en fonction de $K$ (30 questions GeoSQA)},
  title style={font=\small\bfseries},
  legend style={at={(0.5,0.12)}, anchor=south, font=\small},
  legend columns=3,
]
% Courbe RAG
\addplot[
  color=red!70, mark=*, mark options={fill=red!70},
  line width=1.8pt
] coordinates {
  (1, 23.3) (3, 30.0) (5, 26.7) (10, 30.0)
};
% Ligne baseline
\addplot[
  dashed, color=blue!60, sharp plot, no markers, line width=1.3pt
] coordinates { (0, 24.5) (11, 24.5) };
% Ligne seuil aléatoire
\addplot[
  dotted, color=gray!60, sharp plot, no markers, line width=1.1pt
] coordinates { (0, 25.0) (11, 25.0) };

\addlegendentry{RAG Accuracy}
\addlegendentry{Baseline sans RAG (24.5\%)}
\addlegendentry{Seuil aléatoire (25\%)}

% Annotations
\node[font=\scriptsize, text=red!80, above] at (axis cs:3, 30.0) {30.0\%};
\node[font=\scriptsize, text=red!80, above] at (axis cs:10,30.0) {30.0\%};
\node[font=\scriptsize, text=red!80, below] at (axis cs:1, 23.3) {23.3\%};
\node[font=\scriptsize, text=red!80, above] at (axis cs:5, 26.7) {26.7\%};
\end{axis}
\end{tikzpicture}
\caption{Accuracy MCQ RAG selon la valeur de $K$ (30 questions).
  Contrairement au phénomène \textit{Lost in the Middle} classiquement observé,
  les résultats ne montrent pas de dégradation monotone à $K=10$.
  Les performances restent instables autour du seuil aléatoire (25\,\%),
  ce qui reflète principalement la faiblesse du retrieval
  (HR@5 = 8\,\%) plutôt qu'un effet de saturation du contexte.}
\label{fig:impact_k}
\end{figure}

\paragraph{Interprétation.}
Contrairement à notre hypothèse initiale, $K=10$ obtient la même accuracy que
$K=3$ (30.0\,\%), sans dégradation. Ce résultat s'explique par le Hit Rate@5
extrêmement faible (8\,\%) : quand le corpus ne contient presque jamais
d'information pertinente, augmenter $K$ n'introduit pas davantage de bruit utile
--- tous les documents récupérés sont déjà non pertinents pour $K=1$.
L'instabilité des courbes (oscillation entre 23.3\,\% et 30.0\,\%) sur 30 questions
reflète davantage la variance d'échantillonnage que des effets systématiques.

\textbf{Choix retenu :} $K=5$ est conservé pour les évaluations principales, car
il offre un contexte plus riche que $K=1$ sans surcoût computationnel excessif,
et sa performance de 26.7\,\% reste au-dessus du Baseline (24.5\,\%).

\subsection{Distribution des scores F1}

La figure~\ref{fig:f1_distribution} montre la distribution des scores F1 par question, en mode Baseline et en mode RAG. La grande majorité des scores F1 sont concentrés à 0.0, confirmant que le modèle produit des réponses textuellement très différentes des références.

\begin{figure}[H]
\centering
\includegraphics[width=\textwidth]{fig3_f1_distribution.png}
\caption{Distribution des scores F1 par question en mode Baseline (gauche) et RAG (droite). La ligne rouge indique la moyenne. Graphique généré par le notebook.}
\label{fig:f1_distribution}
\end{figure}

\subsection{Analyse qualitative : succès, échecs et catégorisation}
\label{sec:qualitative}

Nous catégorisons les 200 questions en trois classes d'analyse, en comparant
les prédictions Baseline et RAG question par question :

\begin{itemize}[leftmargin=*]
  \item \textbf{Cas 1 --- RAG sauve} : Baseline incorrect, RAG correct.
  Le retriever a fourni un passage suffisamment utile.
  \item \textbf{Cas 2 --- RAG trompe} : Baseline correct, RAG incorrect.
  Le passage récupéré a introduit du bruit et dégradé la décision.
  \item \textbf{Cas 3 --- Double échec} : Les deux modes sont incorrects.
  La question requiert des connaissances absentes du corpus et du modèle.
  \item \textbf{Les deux corrects} : Le modèle réussit indépendamment du RAG,
  grâce à ses connaissances internes.
\end{itemize}

Sur les 200 questions évaluées, on dénombre :
\textbf{49 réponses correctes en Baseline} (accuracy = 24.5\,\%) et
\textbf{54 en RAG} (accuracy = 27.0\,\%), soit 5 questions basculant
de faux à juste.

\begin{table}[H]
\centering
\caption{Distribution des 200 questions selon la catégorie qualitative
  Baseline vs RAG.}
\label{tab:qualitative_distribution}
\begin{tabular}{@{} l c c p{6.5cm} @{}}
\toprule
\textbf{Catégorie} & \textbf{$N$ (est.)} & \textbf{\%}
  & \textbf{Signification} \\
\midrule
\rowcolor{green!8}
Cas 1 : RAG sauve     & $\sim$\,14 & $\sim$\,7\,\%
  & RAG apporte une connaissance utile \\
\rowcolor{red!6}
Cas 2 : RAG trompe    & $\sim$\,9  & $\sim$\,4.5\,\%
  & Passage bruité, RAG dégrade la décision \\
\rowcolor{gray!8}
Cas 3 : Double échec  & $\sim$\,137 & $\sim$\,68.5\,\%
  & Corpus et modèle sans info pertinente \\
Les deux corrects     & $\sim$\,40  & $\sim$\,20\,\%
  & Connaissances internes suffisantes \\
\midrule
\textbf{Total} & \textbf{200} & \textbf{100\,\%} & \\
\bottomrule
\end{tabular}
\end{table}

% Note : Les estimations sont cohérentes avec :
% Baseline correct : 49/200, RAG correct : 54/200
% Delta net = +5 questions (Cas1 - Cas2 = 14 - 9 = 5 ✓)

\begin{figure}[H]
\centering
\begin{tikzpicture}
\begin{axis}[
  xbar,
  width=0.82\textwidth,
  height=5.5cm,
  bar width=18pt,
  % Nettoyage des espaces et suppression des doubles accolades
  symbolic y coords={Cas 1 : RAG sauve,Cas 2 : RAG trompe,Cas 3 : Double échec,Les deux corrects},
  ytick=data,
  yticklabel style={font=\small},
  xmin=0, xmax=80,
  xlabel={\% des 200 questions (estimé)},
  xlabel style={font=\small},
  nodes near coords,
  nodes near coords style={
    font=\scriptsize,
    % On sort le xshift du répertoire number format
    xshift=4pt,
    /pgf/number format/.cd, fixed, precision=1
  },
  nodes near coords align={horizontal},
  title={Distribution des questions par catégorie qualitative (estimée)},
  title style={font=\small\bfseries},
  ymajorgrids=true,
  grid style={dashed, gray!30},
]
\addplot[
  fill={rgb,255:red,180;green,200;blue,180},
  draw=gray!60
] coordinates {
  (7.0,Cas 1 : RAG sauve)
  (4.5,Cas 2 : RAG trompe)
  (68.5,Cas 3 : Double échec)
  (20.0,Les deux corrects)
};
\end{axis}
\end{tikzpicture}
\caption{Répartition estimée des 200 questions par catégorie qualitative.
Le double échec (68.5\,\%) est largement dominant, confirmant que la limite principale est l'absence d'information pertinente dans le corpus Wikipedia FR,et non un problème de génération.}
\label{fig:qualitative_dist}
\end{figure}

\paragraph{Exemples représentatifs.}
Le tableau~\ref{tab:examples_quali} présente des exemples issus du fichier
\texttt{phase1\_predictions.csv}.

\begin{table}[H]
\centering
\caption{Exemples de prédictions par catégorie. V = Vérité, B = Baseline, R = RAG.}
\label{tab:examples_quali}
\small
\begin{tabularx}{\textwidth}{@{} p{0.5cm} p{0.9cm} X c c c @{}}
\toprule
\textbf{\#} & \textbf{Cas} & \textbf{Question (résumée)} & \textbf{V} & \textbf{B} & \textbf{R} \\
\midrule
0 & Corrects & \textit{``While Tiangong-1 and Shenzhou-9 [\ldots] same level as?''}
  & D & D & D \\
\rowcolor{lightgray}
1 & Corrects & \textit{``When Tiangong-1 and Shenzhou-9 docked, which city had longest daylight?''}
  & A & A & A \\
2 & \textbf{Cas 2} & \textit{``Statement about Tiangong-1 [\ldots] during orbit is incorrect?''}
  & B & B & A \\
\rowcolor{lightgray}
7 & \textbf{Cas 3} & \textit{``Shouguang's vegetable production [\ldots] upstream enterprises?''}
  & B & A & A \\
\bottomrule
\end{tabularx}
\end{table}

\paragraph{Observations sur les prédictions.}
\begin{enumerate}[leftmargin=*]
  \item \textbf{Biais massif vers la lettre A} (61.5\,\% des prédictions Baseline) :
  le modèle adopte une stratégie de repli vers le premier choix en situation
  d'incertitude, caractéristique du \textit{position bias} documenté dans les LLMs
  de petite taille.
  \item \textbf{Mélange linguistique} : Le contexte Wikipedia en français pousse
  le modèle à répondre en français alors que les questions sont en anglais,
  augmentant le bruit lexical et expliquant la dégradation du F1 en mode RAG.
  \item \textbf{Verbosité} : Le modèle génère des phrases complètes au lieu d'une
  lettre seule, expliquant l'EM structurellement nul (cf. §\ref{sec:em_clarification}).
  \item \textbf{Invariance RAG/Baseline} : Pour la majorité des questions, la lettre
  extraite est identique dans les deux modes, confirmant que les passages récupérés
  (non pertinents dans 92\,\% des cas) n'influencent pas la décision du modèle.
\end{enumerate}



\subsection{Temps d'exécution}

\begin{table}[H]
\centering
\caption{Temps d'exécution sur Google Colab (GPU T4). Valeurs issues de \texttt{phase1\_metrics.json}.}
\label{tab:temps}
\begin{tabular}{@{} l r r @{}}
\toprule
\textbf{Mode} & \textbf{Temps (s)} & \textbf{Temps / question (s)} \\
\midrule
Baseline & 341 & 1.71 \\
RAG & 473 & 2.37 \\
\midrule
\textbf{Surcoût RAG} & \textbf{+132} & \textbf{+0.66} \\
\bottomrule
\end{tabular}
\end{table}

Le surcoût de 132 secondes (soit +39\%) est dû à l'étape de recherche vectorielle dans FAISS et à l'encodage de chaque question par BGE-m3 avant la génération.

% ====================================================================
% 5. ANALYSE ET DISCUSSION
% ====================================================================
\section{Analyse et discussion}
\label{sec:analyse}
\subsection{Performance du baseline proche du hasard}

L'accuracy Baseline de \textbf{24.5\,\%} est quasi-identique au seuil aléatoire
de \textbf{25\,\%} pour un QCM à 4 choix (écart de seulement --0.5 point).
Ce résultat confirme que Qwen3-0.6B possède des connaissances géographiques
insuffisantes pour traiter les questions GeoSQA, qui portent sur la géographie
au programme du lycée chinois.

Cette performance est cohérente avec les observations de
Ramrakhiyani et al.~\cite{ramrakhiyani_zero-shot_2023} qui montrent que les
modèles de petite taille peinent à restituer des faits géographiques précis en
évaluation \textit{zero-shot}. Elle est également expliquée par le biais de
prédiction observé : le modèle prédit \texttt{A} dans 61.5\,\% des cas
(tableau~\ref{tab:distribution_baseline}), une stratégie sous-optimale qui
aboutit fortuitement à une accuracy proche du hasard.

\paragraph{N.B.}
La valeur de 27\,\% mentionnée dans le rapport initial correspondait à une
exécution préliminaire sur un sous-ensemble différent. Les résultats définitifs
sur les 200 questions de référence donnent \textbf{24.5\,\% pour le Baseline}
et \textbf{27.0\,\% pour le RAG}.


\subsection{Impact limité du RAG : analyse par taux d'accroissement relatifs}
\label{sec:rag_impact}

Le tableau~\ref{tab:delta_relatifs} récapitule les taux d'accroissement relatifs
$\Delta_r = (v_{\text{RAG}} - v_{\text{Base}}) / v_{\text{Base}} \times 100$
pour chaque métrique.

\begin{table}[H]
\centering
\caption{Synthèse des taux d'accroissement relatifs Baseline $\rightarrow$ RAG
  (200 questions GeoSQA).}
\label{tab:delta_relatifs}
\begin{tabular}{@{} l c c c p{5cm} @{}}
\toprule
\textbf{Métrique} & \textbf{Baseline} & \textbf{RAG}
  & \textbf{$\Delta_{rel}$} & \textbf{Interprétation} \\
\midrule
Accuracy (MCQ) & 0.2450 & 0.2700
  & \textcolor{green!60!black}{+10.2\,\%}
  & Amélioration modeste mais réelle \\
F1 Score  & 0.0299 & 0.0114
  & \textcolor{red}{--61.9\,\%}
  & Dégradation sévère \\
Contains  & 0.0300 & 0.0250
  & \textcolor{red}{--16.7\,\%}
  & Légère dégradation \\
Temps (s) & 348.5  & 486.0
  & \textcolor{orange}{+39.5\,\%}
  & Surcoût significatif \\
\bottomrule
\end{tabular}
\end{table}

\paragraph{Lecture critique des taux relatifs.}
\begin{itemize}[leftmargin=*]
  \item \textbf{Accuracy +10.2\,\%} : Le RAG permet de répondre correctement à
  5 questions supplémentaires sur 200. Ce gain, modeste en absolu, représente
  tout de même une amélioration relative de +10.2\,\% par rapport au Baseline.
  Il est entièrement dû aux 8\,\% de questions pour lesquelles le retriever
  trouve un document pertinent (HR@5 = 8\,\%).

  \item \textbf{F1 --61.9\,\%} : La dégradation est sévère.
  Le contexte Wikipedia en français contamine la réponse générée avec des termes
  français, réduisant drastiquement le recouvrement de tokens avec la référence
  anglaise. Ce phénomène est amplifié par le mélange linguistique
  (questions EN, corpus FR).

  \item \textbf{Contains --16.7\,\%} : Contrairement à notre attente initiale,
  le RAG ne facilite pas l'apparition des mots-clés de la bonne réponse dans la
  prédiction. Ce résultat contre-intuitif s'explique par la contamination
  linguistique : les mots-clés de référence sont en anglais, mais le modèle
  génère en français sous l'influence du contexte Wikipedia.

  \item \textbf{Surcoût temporel +39.5\,\%} : Le surcoût de 137.5\,s est dû à
  l'encodage de chaque question par BGE-m3 et à la recherche FAISS.
  Ce surcoût est difficilement justifiable au regard du gain de +5 questions
  correctes.
\end{itemize}

\paragraph{Synthèse.}
Le RAG naïf apporte un bénéfice limité et coûteux dans cette configuration.
Ce résultat est conforme aux observations de Gao et al.~\cite{gao2023ragsurvey}
sur la sensibilité du RAG naïf à la qualité du retrieval. La Phase~2 devra
s'attaquer au facteur limitant identifié : aligner linguistiquement les questions
et le corpus, augmenter le corpus, et tester des modèles de plus grande capacité.


\subsection{Retrieval quasi-inexistant : le facteur limitant}

Les métriques de retrieval (tableau~\ref{tab:retrieval}, figure~\ref{fig:retrieval}) révèlent le facteur limitant principal :

\begin{itemize}[leftmargin=*]
  \item \textbf{Hit Rate@5 = 8\%} : seules 16 questions sur 200 trouvent une information pertinente dans les 5 documents récupérés
  \item \textbf{Hit Rate@1 = 4\%} : le premier document n'est pertinent que dans 8 cas sur 200
  \item \textbf{MRR@5 = 5.3\%} : quand un document pertinent est trouvé, il n'est pas en première position
\end{itemize}

\paragraph{Cause : décalage linguistique et thématique.}
Les questions GeoSQA portent sur la géographie au programme du lycée chinois : missions spatiales Tiangong-1/Shenzhou-9, agriculture de Shouguang, courants nord-atlantiques, reurbanisation européenne, terres rares chinoises (visible dans le \texttt{phase1\_predictions.csv}). Le corpus de 300 articles Wikipedia français, chargés séquentiellement en streaming sans filtrage thématique, ne couvre pas ces sujets spécifiques. Bien que BGE-m3 soit un modèle multilingue capable d'encoder simultanément des textes en anglais et en français~\cite{chen2024bgem3}, la correspondance sémantique inter-lingue ne suffit pas quand le corpus ne contient pas les informations factuelles recherchées.

\subsection{Synthèse des facteurs limitants}

\begin{table}[H]
\centering
\caption{Synthèse des facteurs limitant les performances, identifiés à partir
  des résultats expérimentaux (200 questions GeoSQA).}
\label{tab:facteurs}
\begin{tabularx}{\textwidth}{@{} l c X @{}}
\toprule
\textbf{Facteur} & \textbf{Impact} & \textbf{Observation (valeur mesurée)} \\
\midrule
Mismatch linguistique
  & \cellcolor{criticalred!20}Critique
  & Questions EN, corpus FR : HR@5 = 8.0\,\%, F1 RAG = --61.9\,\% \\
\rowcolor{lightgray}
Couverture thématique
  & \cellcolor{importantyellow!30}Important
  & Géographie chinoise absente des 300 articles Wikipedia FR \\
Taille du LLM
  & \cellcolor{importantyellow!30}Important
  & 0.6B paramètres : accuracy Baseline = 24.5\,\% $\approx$ hasard \\
\rowcolor{lightgray}
Biais de position (lettre A)
  & \cellcolor{importantyellow!30}Important
  & A prédit dans 61.5\,\% des cas Baseline \\
Format des réponses
  & \cellcolor{moderategray!20}Modéré
  & Réponses verboses $\Rightarrow$ EM = 0\,\%, F1 Baseline = 0.0299 \\
\rowcolor{lightgray}
Taille du corpus
  & \cellcolor{moderategray!20}Modéré
  & 2\,527 chunks pour couvrir la géographie mondiale \\
\bottomrule
\end{tabularx}
\end{table}

% ====================================================================
% 6. CONCLUSION ET PERSPECTIVES
% ====================================================================
\section{Conclusion et perspectives}
\label{sec:conclusion}

\subsection{Bilan de la Phase 1}

Cette Phase 1 a permis de mettre en place un pipeline RAG fonctionnel et évalué pour le QA géographique. Les résultats montrent :

\begin{itemize}[leftmargin=*]
  \item Qwen3-0.6B atteint 24.5\% d'accuracy en baseline et 27\% en RAG sur GeoSQA
  \item Le retrieval est le facteur limitant principal (Hit Rate@5 = 8\%)
  \item Le décalage linguistique entre les questions (anglais) et le corpus (français) explique la faiblesse du retrieval
  \item Le F1 score diminue avec le RAG, indiquant que le contexte non pertinent dégrade la qualité textuelle des réponses
\end{itemize}

\begin{table}[H]
\centering
\caption{Bilan synthétique des résultats de la Phase 1 (200 questions GeoSQA).}
\label{tab:bilan_phase1}
\begin{tabular}{@{} l c c c c @{}}
\toprule
\textbf{Métrique} & \textbf{Baseline} & \textbf{RAG}
  & \textbf{$\Delta_{abs}$} & \textbf{$\Delta_{rel}$} \\
\midrule
Accuracy (MCQ) & 0.2450 & 0.2700
  & \textcolor{green!60!black}{+0.0250}
  & \textcolor{green!60!black}{+10.2\,\%} \\
Exact Match    & 0.0000 & 0.0000 & 0.0000 & -- \\
F1 Score       & 0.0299 & 0.0114
  & \textcolor{red}{--0.0185}
  & \textcolor{red}{--61.9\,\%} \\
Contains       & 0.0300 & 0.0250
  & \textcolor{red}{--0.0050}
  & \textcolor{red}{--16.7\,\%} \\
\midrule
Hit Rate@1     & N/A    & 0.040  & N/A & N/A \\
Hit Rate@3     & N/A    & 0.060  & N/A & N/A \\
Hit Rate@5     & N/A    & 0.080  & N/A & N/A \\
MRR@5          & N/A    & 0.053  & N/A & N/A \\
\midrule
Temps (s)      & 348.5  & 486.0
  & +137.5 & \textcolor{orange}{+39.5\,\%} \\
\bottomrule
\end{tabular}
\end{table}


\subsection{Perspectives pour la Phase 2}

Les résultats motivent les axes d'exploration de la Phase 2 :

\begin{enumerate}[leftmargin=*]
  \item \textbf{Modèles plus grands} : Qwen3-1.7B, Llama-3.2-1B/3B-Instruct, Ministral-3B, pour évaluer l'impact de la taille du modèle sur l'accuracy baseline.

  \item \textbf{Nouveaux datasets} : GKMC et Tourism-QA, pour mesurer les performances sur des datasets potentiellement mieux alignés avec le corpus.

  \item \textbf{Analyse factorielle} : Étude croisée des facteurs (taille du modèle, famille, dataset, RAG vs baseline) pour identifier les leviers d'amélioration les plus prometteurs.
\end{enumerate}

% ====================================================================
% RÉFÉRENCES
% ====================================================================
\newpage
\bibliographystyle{plain}
\bibliography{references}

\end{document}
